% ============================================================
% Chapter 6: Conclusion
% ============================================================
\chapter{Conclusion}
\label{ch:conclusion}

\section{Summary of Contributions}
\label{sec:summary}

This thesis has presented the Adaptive Subspace Bayesian Optimisation (ASBO)
framework, a principled approach to scaling Bayesian optimisation to
high-dimensional problems.  The key contributions are:

\begin{enumerate}
    \item \textbf{Adaptive subspace identification:} We developed a variational
    inference procedure that learns the intrinsic low-dimensional structure of
    the objective function, adapting to varying local dimensionality across the
    search space.  This eliminates the rigid fixed-subspace assumption of prior
    methods.

    \item \textbf{Information-theoretic acquisition:} We derived the Adaptive
    Entropy Search (AES) acquisition function, which naturally adapts to the
    local dimensionality and provides stronger exploration guarantees than
    \gls{ucb}-based alternatives.

    \item \textbf{Theoretical guarantees:} We established convergence bounds
    for ASBO (\cref{thm:regret}), showing that the simple regret depends on the
    intrinsic dimensionality rather than the ambient dimension.  This provides a
    formal justification for the dimensionality reduction approach.

    \item \textbf{Empirical validation:} We demonstrated state-of-the-art
    performance on 12 synthetic benchmarks and 3 real-world engineering
    problems, with up to $3.2\times$ improvement over the best baseline.

    \item \textbf{Open-source library:} We released the \texttt{highdimbo}
    Python library, enabling practitioners to apply ASBO to their own
    optimisation problems with minimal configuration.
\end{enumerate}

\section{Limitations}
\label{sec:limitations}

Despite the strong results, several limitations of the current framework should
be acknowledged:

\begin{itemize}
    \item \textbf{Computational overhead:} The adaptive subspace inference adds
    computational cost compared to simpler methods like REMBO.  While the
    overhead is modest for $d < 200$, it becomes more significant for very high
    dimensions.

    \item \textbf{Linear subspace assumption:} The current framework assumes a
    linear subspace structure.  Extending ASBO to non-linear manifold
    structures (e.g., using deep \glspl{vae}) is an important direction for
    future work.

    \item \textbf{Parallel evaluations:} The current implementation evaluates
    one point at a time.  Extending ASBO to batch acquisition (parallel
    evaluations) would significantly improve wall-clock efficiency.

    \item \textbf{Constraint handling:} The framework currently addresses
    unconstrained optimisation.  Extending ASBO to handle black-box constraints
    (e.g., via constraint acquisition functions) would broaden its applicability.
\end{itemize}

\section{Future Work}
\label{sec:future-work}

We identify several promising directions for future research:

\begin{enumerate}
    \item \textbf{Non-linear subspace models:} Replace the linear projection
    with deep generative models (e.g., \glspl{vae} or normalising flows) to
    capture non-linear manifold structures in the objective function.

    \item \textbf{Batch ASBO:} Develop batch acquisition strategies that select
    multiple evaluation points simultaneously, leveraging recent advances in
    batch \gls{bo}.

    \item \textbf{Multi-fidelity extension:} Combine ASBO with multi-fidelity
    optimisation, using cheap low-fidelity evaluations to inform the subspace
    estimate before committing to expensive high-fidelity evaluations.

    \item \textbf{Constraint handling:} Integrate constraint acquisition
    functions to handle black-box equality and inequality constraints.

    \item \textbf{Theoretical extensions:} Tighten the convergence bounds in
    \cref{thm:regret} and extend the analysis to non-Gaussian likelihoods.

    \item \textbf{Applications:} Apply ASBO to additional domains, including
    reinforcement learning hyperparameter tuning, materials discovery, and
    climate model calibration.
\end{enumerate}

\section{Broader Impact}
\label{sec:broader-impact}

The work presented in this thesis has potential broader impact across several
domains.  In drug discovery, ASBO can accelerate the identification of
promising drug candidates by reducing the number of expensive wet-lab
experiments required, potentially shortening the timeline from target
identification to clinical trials.  In materials science, the framework can
enable more efficient exploration of compositional spaces, accelerating the
discovery of novel materials with tailored properties.  In machine learning,
ASBO can reduce the computational cost of hyperparameter tuning, making
advanced model architectures accessible to researchers with limited
computational resources.

We note that the open-source release of the \texttt{highdimbo} library lowers
the barrier to adoption, enabling broad access to these optimisation
capabilities.  We encourage the community to build upon this work and explore
applications in their respective domains.

\section{Lessons Learned}
\label{sec:lessons}

Several practical lessons emerged during this research that may benefit future
work:

\begin{enumerate}
    \item \textbf{Start with the subspace:} Estimating the intrinsic
    dimensionality early in the optimisation process provides a strong inductive
    bias that accelerates convergence, even when the estimate is imperfect.

    \item \textbf{Balance exploration and computation:} Information-theoretic
    acquisition functions provide stronger theoretical guarantees but are more
    computationally expensive than \gls{ucb}-based alternatives.  The choice
    should be guided by the relative cost of function evaluations versus
    acquisition function computation.

    \item \textbf{Validate incrementally:} Each component of the ASBO framework
   ---subspace identification, local surrogate modelling, and acquisition
    function design---was validated independently before integration.  This
    incremental approach facilitated debugging and provided clearer insights
    into the contribution of each component.
\end{enumerate}

\section{Closing Remarks}
\label{sec:closing}

High-dimensional black-box optimisation remains a fundamental challenge across
science and engineering.  By combining adaptive subspace identification with
information-theoretic acquisition functions, ASBO provides a principled and
effective approach that scales to hundreds of dimensions while maintaining
sample efficiency.  We hope that the theoretical insights, algorithms, and
open-source tools presented in this thesis will empower researchers and
practitioners to tackle optimisation problems that were previously intractable.

The journey from low-dimensional Bayesian optimisation to the scalable
framework presented here has been both intellectually rewarding and practically
impactful.  As the complexity of the problems we face continues to grow, we are
confident that principled, adaptive approaches like ASBO will play an
increasingly important role in scientific discovery and engineering innovation.
